\section{AegisKV Design}
\label{sec:design}

AegisKV manages reusable KV checkpoints across shared DRAM and a larger
SSD or shared-storage tier. GPU HBM holds ordinary execution state and
bounded recovery workspace. The protected failure domain is a fail-stop
compute worker; the directory, token log, and advertised backing copies
must survive that failure. A host-local SSD does not protect against loss
of the same host. The following is a design specification; the new
mechanism figures are constructed examples or analytical schedules, not
measurements of an implemented controller.

\subsection{Design Insights}
\label{sec:design-insights}

\paragraph{Prioritize workflow exposure, while preserving normal reuse.}
A faster restore advances a workflow only when its delay exceeds the
available dependency slack (Figure~\ref{fig:motivation-recovery-exposure}).
However, a large-slack agent can reuse its history after a tool call.
Assigning that agent permanently to recomputation would waste serving
capacity. AegisKV therefore separates \emph{preserving a reusable copy}
from \emph{retaining it in the faster tier}. Shared immutable objects
serve both purposes without duplicating their physical byte charge.

\paragraph{Sparse access identifies candidates, not a complete backup.}
Native sparse attention can use only a subset of historical KV for a
query, but future selections remain uncertain. Figure~\ref{fig:design-insights}
shows that causal access history predicts future block demand unevenly
across models. At a 40\% block budget, past-frequency coverage is 77.8\%
for NOSA, 53.6\% for MiniCPM4, and 43.0\% for the experimental Seer path.
The corresponding same-budget oracle means are 93.7\%, 59.0\%, and
43.3\%. Seer's low absolute coverage is therefore largely a budget limit
for this future-query union, rather than a large gap to its oracle.
These observations motivate selective fast-tier retention and a mandatory
miss path, rather than discarding the remaining checkpoint. The three
documents per model support this observed contrast, not a population-wide
prediction guarantee or a measured recovery speedup.

\begin{figure}[t]
 \centering
 \includegraphics[width=\columnwidth]{figures/retry-kv-prediction.pdf}
 \caption{Historical access predicts distinct future KV unevenly.
 Three native PG19 captures/model use 32K history: queries
 32,257--32,768 predict blocks needed by 32,769--32,832.
 Objects are layer/64-token blocks containing all KV heads;
 mandatory blocks count against the 40\% budget. Seer uses an
 experimental Qwen3-8B Prefill gate~\cite{song2026seerqwen3}, threshold
 $10^{-3}$, with final-block dense override disabled.
 Bars/whiskers are document means/ranges, not recovery runs. Diamonds
 show offline oracle mean coverage with the same mandatory blocks and
 capacity; future demand is unavailable to the online predictor.}
 \label{fig:design-insights}
 \Description{Held-out native block coverage varies across three models;
 weak effects remain visible.}
\end{figure}

\paragraph{Overlap determines the marginal value of fast memory.}
DRAM-resident KV avoids the SSD read but still requires H2D transfer.
Early-layer reads can delay the start of replay, whereas later reads may
overlap with preceding replay computation. Their values also depend on
the current checkpoint length, uncommitted tail, demand reveal time,
resource contention, and workspace. A fixed resident layer count cannot
represent all requests. AegisKV evaluates feasible layer/block plans and
uses online rent to screen competing DRAM claims.

\subsection{System Overview}
\label{sec:system-overview}

Figure~\ref{fig:system-overview} organizes AegisKV into three layers.
The \emph{Application} layer issues agent requests and exposes workflow
dependencies and progress. In the \emph{Controller} layer, the Tracker
estimates remaining durations and slack, the KV Master values placements
and maintains rental credits, and the Recovery-aware Router dispatches
requests and selects surviving recovery targets. A surviving directory
records immutable objects, location generations, and checkpoint manifests,
alongside the token log. Within the Master, layer-I/O profiles feed the
workflow cost model; independent DRAM/storage credit ledgers screen
competitors; a joint DAG check and resource reservation authorize the
complete plan. The directory supplies manifest versions and physical
read leases for this check. The \emph{Engine} layer contains interchangeable
GPU workers, their HBM working sets, and the shared DRAM/storage data path.
Engine hooks report sealed state and resource pressure, while transfer
queues apply the Master's decisions. Shared backing resides outside the
protected compute-worker failure boundary, allowing a compatible healthy
DP to restore state without a preassigned peer-memory target.

\begin{figure}[t]
 \centering
 \includegraphics[width=\columnwidth]{figures/system-overview.pdf}
 \caption{AegisKV's three-layer architecture. Controller exposes its
 prediction, valuation, competition, and commit checks. D/S in Rent
 denote separate DRAM/storage credit ledgers; the joint DAG check and
 reservation authorize a placement. Directory state includes manifests,
 token logs, and leases. Agents A--D illustrate a fork/join workflow.
 Blue bidirectional links carry KV; dark dashed links carry
 control/metadata; red dashes illustrate retry. Shared DRAM and
 SSD/storage sit outside the DP worker boundaries.
 HBM outlines mark demand candidates, and the last column is the unsealed
 tail. GPU symbols abstract each DP's GPU/TP group.}
 \label{fig:system-overview}
 \Description{Application, Controller, and Engine bands. Tracker applies
 GBDT and critical-path analysis. KV Master combines layer-I/O and WCT
 cost, tier-specific rent, joint DAG revaluation and reservation. Router
 uses load and locality and fences failures. Directory versions and
 leases guard placement. Symmetric DP GPU/HBM workers share backing.}
\end{figure}

A physical object contains one layer's KV shard for a sealed token block.
Identity includes model/weight revision, prefix and position semantics,
layer, TP shard, format, and token interval. A legal placement group
includes the objects and selector state required by the native engine;
logical sparse selections, physical storage blocks, and kernel tiles need
not coincide. The engine chooses the physical block size; the 64-token
unit in Figure~\ref{fig:design-insights} is an experimental convention.

Prefill offers newly sealed groups. During decode, a configured,
block-aligned interval $Z$ triggers backup review. Arrivals, tool handoffs,
material demand/slack changes, and capacity pressure also trigger review.
A decision does not imply immediate commitment: if $M$ is the last
complete checkpoint frontier, the replay tail $N$ may exceed $Z$ when
copies lag. A checkpoint includes all historical dependencies needed to
resume native attention; DRAM and storage may jointly hold its objects.

\subsection{Workflow Tracker: Remaining Time and Slack}
\label{sec:workflow-protection}

The Tracker builds a residual DAG of unfinished execution segments.
Fork/join boundaries split parent work so child execution is not counted
again as parent waiting. A gradient-boosted decision-tree model (GBDT)
predicts remaining duration. Its inputs group three types of causal
information: role/model and current input length; elapsed execution and
completed-call count; and completed calls' output lengths, API durations,
context growth, and inter-call gaps. Tool type and task descriptors can
be added when the orchestrator exposes them. Current-call output length,
future tool results, and future call count are unavailable inputs.
Entry predictions mask missing history; elapsed time is not subtracted
again from an already remaining-time prediction.

Figure~\ref{fig:tracker-correlations} examines six observable numerical
inputs in the CC/Weka trace~\cite{semianalysis2026cctraces}. Each of 1,443
eligible agents contributes one snapshot at its second request, after
its first call completes. The target is the remaining duration of the
whole agent, including later calls and gaps. Input-length features are
64-token hash-prefix proxies; explicit task semantics and tool outcomes
are not recorded. All six marginal correlations are weak, and their
workflow-bootstrap intervals cross zero. Thus, a single-feature rule
is not established by this cohort. These associations neither measure
GBDT feature importance nor validate its nonlinear predictions or the
residual-segment target required by the online Tracker.

\begin{figure}[t]
 \centering
 \includegraphics[width=\columnwidth]{figures/tracker-correlations.pdf}
 \caption{Tracker evidence at request 2: 1,443 agents in 170 workflows,
 one snapshot per agent. (a) Spearman associations with whole-agent
 remaining time; bars are percentile 95\% intervals from 2,000
 whole-workflow bootstrap samples. Asterisks mark prefix-block length
 proxies. (b) Empirical absolute-error CDFs from existing five-fold,
 workflow-disjoint predictions: training-fold median, basic GBDT, and
 GBDT with call history. All tails are retained; the horizontal scale is
 logarithmic above 1\,s. Neither panel validates residual-segment Slack.}
 \label{fig:tracker-correlations}
 \Description{Six weak feature associations and whole-agent held-out
 error distributions. GBDT errors remain substantial; history does not
 uniformly improve the basic model.}
\end{figure}

Figure~\ref{fig:tracker-correlations}(b) makes prediction error explicit.
The existing workflow-separated five-fold pilot gives mean absolute
errors of 438.6, 406.7, and 407.1\,s for median, basic GBDT, and
history GBDT. Their median absolute percentage errors are 76.2\%,
67.1\%, and 69.3\%. The basic-minus-median MAE difference has a paired
workflow-bootstrap 95\% interval of $[-80.9,4.9]$\,s; it does not
establish a reliable mean-error reduction. Fitted predictions are fixed
in this bootstrap. Only 36.5\% of basic-GBDT predictions are within 50\%
of actual remaining time. The cohort comprises completed subagents and
includes future inter-call gaps, so it cannot validate runtime segment
predictions. Deployment validation must separate whole workflows and
time periods and report false-positive slack, not just duration error.
Published trajectory
prediction~\cite{zhang2026heddle} motivates feature construction but does
not establish wall-clock accuracy for AegisKV.

Critical-path analysis~\cite{kelley1959cpm} computes earliest and latest
finishes; their nonnegative difference is slack $\sigma_r$ inherited by a
request's segment. Requests on the same segment do not each receive an
independent copy of that margin. A refresh costs $O(|V|+|E|)$ after
prediction. Missing dependencies or invalid predictions fall back to
request recovery cost and normal reuse with zero slack. A new DAG version
invalidates outstanding valuation snapshots.
Without validated segment forecasts, the policy uses the zero-slack
fallback rather than treating the pilot's point estimates as safe
demotion margins. The current Tracker evidence motivates this guard;
it does not establish predictive accuracy as a system contribution.



\subsection{Recovery-aware Router}
\label{sec:recovery-router}

Normal dispatch balances load and compatible prefix locality, retaining
Agent affinity when useful. The CC/Weka trace~\cite{semianalysis2026cctraces}
shows 97.90\% potential same-agent token coverage, with 0.37\% additional
coverage from other agents; these are potential prefixes, not cache hits. A tool handoff drops
an execution reference but does not delete the cached history or withdraw
its recovery claim.

\begin{figure}[t]
 \centering
 \includegraphics[width=\columnwidth]{figures/agent-prefix-reuse.pdf}
 \caption{Agent-local prefix reuse in the CC/Weka trace:
 98,827 requests, 2,090 agents, and 393 workflows. Coverage uses earlier
 completed requests with the same recorded model and workflow, measured
 by prefix-block hashes. Same-agent and other-agent coverage overlap;
 only the additional cross-agent portion extends same-agent coverage.
 These are potential token-prefix reuse rates, not runtime cache hits.}
 \label{fig:agent-prefix-reuse}
 \Description{Same-agent history covers most observed input tokens;
 other-agent history adds little overall but helps first subagent calls.}
\end{figure}

Figure~\ref{fig:agent-prefix-reuse} separates overlapping coverage from
additional coverage. Cross-agent reuse can still matter on a subagent's
first call; affinity is therefore a preference, not a prohibition on
sharing or load balancing.

On a confirmed failure, the Router fences the old source epoch, resolves
the last complete manifest, and chooses a healthy, layout-compatible
replica with feasible workspace. Queuing and readmission remain part of
recovery cost. Before acknowledging output, a surviving ordered log
records request ID, token position/value, and source epoch; prompt and
accepted tool inputs are also recoverable. Replay suppresses duplicate
delivery and rejects stale callbacks. Exactly-once external tool side
effects remain the orchestrator's responsibility.

\subsection{KV Master: Sparse Recovery Value and Online Rent}
\label{sec:kv-master}

\subsubsection{Valuing retention and demotion}
\label{sec:placement-value}

The Master compares complete, achievable placement plans. For an isolated
interrupted segment, recovery time $R_r$ exposes delay
$\max(0,R_r-\sigma_r)$. A common worker failure can interrupt several
branches: if two required branches each finish 10\,s late, accelerating
only one to 2\,s leaves their join 10\,s late. Thus, isolated block scores
cannot be added to obtain workflow benefit.

For current plan $P$ and candidate $P'$, we evaluate
\begin{equation}
 \Delta J=A(P',P)+\sum_e p_e\sum_w
       \bigl[T_w(P',e)-T_w(P,e)\bigr].
 \label{eq:value}
\end{equation}
Episodes $e$ cover normal next use and named compute-domain failures over
a fixed horizon. The weights $p_e$ are declared probabilities or expected
episode counts, not raw failure rates. A failure includes every request
interrupted in that domain. Profiles estimate additional segment delays
from surviving coverage, tail replay, I/O and target queues; a max-plus
DAG pass yields $T_w$. The transition term $A$ accounts for foreground
delay from copying and temporary occupancy, excluding costs already in
the episode terms. All terms have units of time.

Normal reuse remains in this objective even when failure probability is
zero or slack is large. DRAM demotion is valued against its actual
storage-backed continuation; last-copy deletion is valued against a
remaining checkpoint or reconstruction. Every shared owner is considered,
while each physical replica consumes capacity once. Pending copies help
only in episodes for which the profile predicts successful commitment.
No future failure time, prompt, or query is available to the policy.

\paragraph{Why a priority label is insufficient.}
Slack measures tolerated delay, not how much delay a placement can remove.
Figure~\ref{fig:retention-value} applies the isolated-fault exposure rule
to two declared recovery profiles. At zero slack, retaining A saves
2\,s using two capacity units; at 6\,s slack, retaining B saves 10\,s
using one. Sorting only by slack reverses their value-per-byte order.
Recency can likewise miss a valuable checkpoint with no recent normal
access, while raw saved time ignores how many bytes must remain resident.
These observations require cost- and size-aware comparison. They do not
by themselves select a particular replacement algorithm.

\begin{figure}[t]
 \centering
 \includegraphics[width=\columnwidth]{figures/retention-value.pdf}
 \caption{An analytical counterexample to slack-only priority.
 Per-unit saving is $([R_{\mathrm{slow}}-\sigma]^+
 -[R_{\mathrm{fast}}-\sigma]^+)/b$, conditional on an isolated failure.
 A occupies two normalized capacity units with restore times 4/2\,s;
 B occupies one with times 18/8\,s. Marked points compare A at slack 0
 and B at slack 6. Normal-use value is set to zero to isolate the effect.
 Inputs are constructed, not hardware measurements or evidence that
 rent outperforms direct revaluation.}
 \label{fig:retention-value}
 \Description{A zero-slack group can have much lower saved workflow
 time per unit of capacity than a group with positive slack.}
\end{figure}

\subsubsection{Sparse candidates and layer overlap}
\label{sec:sparse-placement}

For each request and layer, a bounded access summary ranks historical
blocks and includes native mandatory blocks. Candidates comprise legal
individual groups and small layer-ordered bundles, limited by a configured
work budget. The retained count $K_{r,l}$ describes the actual committed
DRAM groups; it can change as context, access, and competitors change.
There is no common resident depth $L$ or fixed attention-selection budget.
A sparse adapter predicts data demand without changing the model's native
selection. A miss reads storage or follows the declared reconstruction
fallback, as Figure~\ref{fig:sparse-recovery} illustrates.

\begin{figure}[t]
 \centering
 \includegraphics[width=\columnwidth]{figures/sparse-recovery.pdf}
 \caption{Sparse retention preserves full recovery semantics.
 (a) A complete historical checkpoint spans DRAM and SSD; dashed tail
 blocks still require replay. (b) For layer 1, a constructed prediction selects
 B0/B3, whereas the native query needs B0/B2. B2 is fetched on demand;
 B3 is an unused prefetch. Blue and tan consistently denote physical
 DRAM and SSD residency (D/S); unselected entries are unfilled.}
 \label{fig:sparse-recovery}
 \Description{Three layers have a complete mixed-tier history and an
 uncommitted tail. A predicted-block mismatch triggers storage access.}
\end{figure}

The profile schedules SSD reads, H2D copies, and layer computation with
finite buffers. DRAM retention removes only the SSD stage. Only remaining
replay computation can hide transfers; avoided historical prefill is not
available overlap time. Figure~\ref{fig:design-pipeline} shows why equal
bytes at different layers can have different value. Figure~\ref{fig:overlap-frontier}
shows diminishing returns within each fixed compute profile. Larger
compute cost hides more later I/O but can increase total recovery time.
These idealized schedules establish a possible layer-position effect;
they do not establish its magnitude with predictor misses or contention.
The runtime evaluates the current $M$, a declared distribution or bound
for unknown tail $N$, and demand-reveal constraints. It does not infer a
universal relationship between $Z$ and resident depth.

\begin{figure}[t]
 \centering
 \includegraphics[width=\columnwidth]{figures/recovery-pipeline.pdf}
 \caption{Analytical recovery schedules with equal-sized demand groups
 across four layers: SSD read 4\,ms/group, H2D 1\,ms/group, and replay
 3\,ms/layer. All four groups are correctly predicted at time zero;
 separate transient workspace holds them, and the three resources are
 independent. Total recovery is 20\,ms with all groups on SSD, 19\,ms
 with layer 4 retained in DRAM, and 16\,ms with layer 1 retained.
 These illustrative rates are not hardware measurements.}
 \label{fig:design-pipeline}
 \Description{Three resource timelines compare all-storage, late-layer
 DRAM retention, and early-layer DRAM retention under the same byte budget.}
\end{figure}

\begin{figure}[t]
 \centering
 \includegraphics[width=\columnwidth]{figures/overlap-frontier.pdf}
 \caption{Analytical marginal value of retaining early-layer groups.
 Each curve fixes replay cost and shares the assumptions of
 Figure~\ref{fig:design-pipeline}; its zero-retention point is its matched
 baseline. Flat regions show reads already hidden by replay. The
 retained depth is an example candidate family, not a system-wide quota.}
 \label{fig:overlap-frontier}
 \Description{Recovery latency saturates at different retained group
 counts for three layer compute costs.}
\end{figure}

To capture limited complementarity, the candidate budget also permits
pairs of groups on required branches exposed to the same fault. Every
candidate is evaluated with its whole victim set through
Equation~\ref{eq:value}. Bounded search can still miss beneficial bundles;
we claim neither optimal placement nor submodularity.

\subsubsection{Online admission, rental, and eviction}
\label{sec:placement-plan}

\paragraph{Why a pressure-aged credit ledger?}
The online problem adds a second requirement to the marginal-value
comparison: contexts and candidate sets grow, future reuse is unknown,
and a claim made during one critical phase should not occupy capacity
indefinitely after that phase. Re-evaluating every resident group through
the full DAG at every checkpoint can require a DAG/profile evaluation
per group. An incremental ordering is useful if it reduces that work
without discarding valuable plans.

Landlord~\cite{young2002landlord} and GreedyDual-Size~\cite{cao1997greedydual}
provide a cost/size-aware aging mechanism for online caching. We adapt
that mechanism as a \emph{candidate screen}: a causal use establishes
credit, admitted capacity competition spends it in proportion to occupied
bytes, and new demand can renew it. No pressure means no rent advance;
an agent's tool wait alone does not trigger eviction. Thus, the ledger
combines cost, size and renewal without choosing a universal wall-clock
TTL. Migration hysteresis has a separate role: paying for transitions
and rejecting small predicted gains.

This rationale is an implementation hypothesis, not an inherited
optimality result. Direct current-value revaluation is a necessary
baseline under the same candidates, safety checks, migration costs and
work budget. The current isolated comparison in
Section~\ref{sec:eval-rental}, with costs in Figure~\ref{fig:eval-rental-cost},
does not show a screening-work reduction and has higher total cost with
rent. Accordingly, we do not claim that rent is necessary or that it
currently improves AegisKV. Its justification as a deployed policy
requires decision-work and realized-cost evidence together.

\paragraph{Tier-local accounting.}
Competitors are screened independently in each capacity tier. For a physical replica
$x$ with bytes $b_x$, that tier's rent level $\rho$ and expiration key
$H_x$ represent remaining credit:
\begin{equation}
 c_x=b_x\max(0,H_x-\rho).
 \label{eq:rent}
\end{equation}
Credits are estimated delay saved, and rent has units of delay per byte;
it is not elapsed wall time. Raising rent by $\delta$ transforms an
unrenewed credit to $\max(0,c_x-b_x\delta)$: larger replicas spend the
same credit faster, and expiration orders their remaining credit per
byte. A shared rent offset represents this aging without decrementing
every replica eagerly. This accounting identity supplies no bound on
realized WCT. For each owner, the local proxy adds the next-use probability times
saved restore wait to fault-episode probabilities times the decrease in
$\max(0,R_r-\sigma_r)$, comparing the same group before and after a legal tier removal.
For DRAM this normally means storage access; for a last storage copy it
means the remaining checkpoint plus reconstruction. Each physical replica
has one byte charge and one tier-local ledger.
It only ranks candidates:
shared failures and migration costs still require the joint plan check.
No classical competitive guarantee transfers to this non-additive setting.
Unlike the classical fixed-cost setting, values here change with slack
and overlap, failures couple owners, admission is optional, and physical
leases restrict release. The joint DAG check is therefore authoritative;
stale credit can still cause the screen to miss a profitable candidate.

Credit is divided into owner shares while bytes are charged once. On a
new causal demand/checkpoint event, the affected owner's share is renewed
from its current proxy. Accrued rent is first settled proportionally across
old shares; other owners are not renewed. Repeated events with the same
version do nothing. Value-only refreshes may reduce credit but do not
refund paid rent. Thus, one active agent cannot perpetually renew an
unrelated owner's claim through a shared physical object.

All committed replicas in a tier age with its $\rho$, including replicas
outside its safe victim set $Q$. Membership in $Q$ requires no outstanding
physical read lease or copy dependency and an executable continuation:
a committed or reservable lower copy for demotion, or an explicitly
revised recovery frontier for last-copy removal. A running request's
logical ownership does not permanently pin every physical location. Credit
expiration alone never releases memory (Figure~\ref{fig:rental-competition}).

\begin{figure}[t]
 \centering
 \includegraphics[width=\columnwidth]{figures/rental-competition.pdf}
 \caption{Constructed rental step with equal-size groups. Trial rent
 increases by two credit units per byte, reducing credits 6/2/5 to
 4/0/3. The leased group ages from one to zero but is excluded from the
 safe victim set. The eligible zero-credit group is only a proposed
 demotion: the complete replacement must pass the workflow-value and
 reservation checks before any committed ledger changes. This illustrates
 accounting and safety, not replacement-policy superiority.}
 \label{fig:rental-competition}
 \Description{All DRAM groups pay rent, while a reader lease prevents
 eviction irrespective of its credit.}
\end{figure}

For optional admission, the Master simulates rent advancement using a
private snapshot. An incoming group's key is its credit per byte plus
current rent. While space is insufficient, it considers the eligible
resident with the smallest expiration key. Only an incoming key strictly
above both that key and trial rent can displace it; the trial rent advances
to that expiration and the victim joins the proposed plan. Ties retain the
resident. If no feasible victim remains, the admission fails. A free-space
admission does not advance rent.

The complete candidate-plus-victim plan must have a favorable upper cost
bound in Equation~\ref{eq:value}, with a configured hysteresis margin.
Rejected previews leave rent and locations unchanged. Accepted previews
atomically recheck versions, reserve resources, commit the new rent, and
enqueue transfers. Actual allocation waits for actual released bytes. A plan that needs a
lower-tier victim includes that victim and its recovery consequences in
the same preview; a DRAM action cannot assume unlimited SSD space.
Once execution starts, a later copy failure triggers reconciliation;
it cannot undo already completed migrations by rolling back a preview.

Required capacity relief has no incoming bidder: bounded rental victim
plans compete by predicted loss per releasable byte, subject to safety.
Storage deletion first removes unowned or superseded state. Under further
pressure, the Master may explicitly withdraw an active checkpoint claim,
valuing the resulting reconstruction for every owner. Before deleting its
last replica, it publishes a fallback to a complete older manifest, or to
prompt and surviving token-log replay if no older checkpoint remains.
Readers of the old generation must drain. It cannot retain the old
recoverability claim after deletion. If no valid fallback or executable
victim exists, it rejects/coalesces new writes and preserves the current
frontier. Identical storage capacity and fallback rules apply to baselines;
the ranking policy is the ablation variable.

\par\medskip\noindent\begin{minipage}{\columnwidth}
\raggedright
 \refstepcounter{designalgorithm}\label{alg:placement}
 \textbf{Algorithm \thedesignalgorithm: Online retention with trial rent}
 \par\smallskip
 \begin{tabularx}{\columnwidth}{@{}r@{\quad}X@{}}
 \toprule
 1 & Snapshot DAG, epochs, copies, leases, and byte ledgers \\
 2 & Reuse existing objects; offer newly sealed dependencies \\
 3 & Update causal demand summaries and affected credit shares \\
 4 & Form bounded sparse/layer candidates and safe victims \\
 5 & For each candidate, privately simulate rental displacement \\
 6 & Evaluate the complete feasible plan through the workflow DAG \\
 7 & Reject uncertain or unprofitable optional plans \\
 8 & Recheck versions; reserve destinations and transfer space \\
 9 & Commit trial rent; enqueue the selected placement changes \\
 10 & On pressure, choose safe relief; block admission if none \\
 11 & On completion/failure, reconcile actual copies and capacity \\
 \bottomrule
 \end{tabularx}
 \par\smallskip\noindent\textit{Preview capacity is not free memory; prediction is not a copy.}
 \end{minipage}\par\medskip

Profiles carry model/layout, size, tail, topology, queue load, support
range, uncertainty, and expiry. Missing or stale profiles freeze optional
promotions and fall back to bounded LRU with the same safety protocol.
Heap maintenance can be logarithmic, but complete cost includes candidate
construction, transfer profiling, and DAG passes. The implementation must
report decision latency and the effect of its candidate-work budget.

\paragraph{Applying a placement decision.}
Engine hooks report actual tier capacity, sealed blocks, access summaries,
and execution workspace. The transfer path enforces
reservations and bounded transfer queues for foreground reuse, demanded
recovery, and checkpoint/writeback. Configured service shares and bounded
quanta limit interference; demanded misses precede speculative reads
within the recovery share. These rules do not promise zero serving cost.

\begin{figure}[t]
 \centering
 \includegraphics[width=\columnwidth]{figures/checkpoint-lifecycle.pdf}
 \caption{Checkpoint and location transitions. (a) Pending tokens
 129--192 leave frontier $M=128$ unchanged; only dependency-complete
 commitment advances it to 192. (b) Demotion copies to storage,
 validates and commits the new location, redirects new readers, drains
 old leases, and then frees DRAM. (c) Last-copy deletion first publishes a
 complete older checkpoint or prompt/token-log replay. Numbers are illustrative;
 the diagram specifies ordering, not a completed fault-injection test.}
 \label{fig:checkpoint-lifecycle}
 \Description{A complete manifest frontier differs from an in-flight
 copy; storage commitment and reader drainage precede DRAM release; last-copy removal explicitly changes the recovery path.}
\end{figure}

Objects move through pending, committed, and retiring generations
(Figure~\ref{fig:checkpoint-lifecycle}). Advancing a frontier requires every
layer, TP shard, selector state, token interval, and matching surviving
log prefix. It is not the largest individually copied token position.
DRAM demotion follows copy--verify--commit--redirect--drain--free. The
backend must meet the claimed fault-domain persistence contract; buffered
write completion alone does not imply durable storage. Copy failure keeps
the previous authoritative location and releases only failed reservations.
Duplicate or stale-epoch callbacks cannot advance coverage.

Recovery acquires a manifest and leases its locations. The transfer path
prefetches predicted native groups in layer order into bounded transient
workspace, while the engine replays the acknowledged tail. On-demand
selection resolves misses before executing the affected operation. Dense
layers use full demand; unavailable sparse metadata is restored or
reconstructed before selection. TP ranks obey matching collectives, and
batched replay is used only when it preserves the engine's native
semantics. A failure during prefill also starts from the last complete frontier,
possibly zero. Speculation may improve timing but never defines correctness.
